\documentclass[12pt,a4paper]{article}

\usepackage{amsmath, amsthm, amssymb}
\usepackage{hyperref}
\usepackage{geometry}
\geometry{margin=2.5cm}

\hypersetup{
  pdftitle={Uniformity Requirements for Weil Fingerprint Energies: Obstructions to the Lipschitz Route to [U]},
  pdfauthor={J\'er\^ome Beau},
  pdfkeywords={Weil fingerprint; finite Heisenberg representation; central character; operator perturbation;
    spectral universality; equicontinuity; Gram-Schmidt rank increment; withdrawn claim; Cosmochrony},
  pdfsubject={Cosmochrony spectral geometry programme, paper U1}
}

\newtheorem{theorem}{Theorem}[section]
\newtheorem{lemma}[theorem]{Lemma}
\newtheorem{proposition}[theorem]{Proposition}
\newtheorem{corollary}[theorem]{Corollary}
\theoremstyle{definition}
\newtheorem{definition}[theorem]{Definition}
\theoremstyle{remark}
\newtheorem{remark}[theorem]{Remark}

\DeclareMathOperator{\Heis}{Heis}
\DeclareMathOperator{\Sym}{Sym}
\newcommand{\Cq}{\mathcal{C}_q}
\newcommand{\Gq}{G_q}
\newcommand{\Piq}{\Pi_q}
\newcommand{\cchi}{c_{\mathrm{BI}}}
\newcommand{\spair}{\sigma_{\mathrm{pair}}}
\newcommand{\dpair}{\delta_{\mathrm{pair}}}
\newcommand{\rhoc}{\rho_{q,c}}
\newcommand{\sigc}{\sigma_c}

\bibliographystyle{plain}
\usepackage{doi}

\begin{document}

    \title{%
      Uniformity Requirements for Weil Fingerprint Energies:\\
      Obstructions to the Lipschitz Route to \textup{[U]}\\[0.4em]
      {\large U1 of the Cosmochrony Spectral Geometry Programme}%
    }
    \author{J\'er\^ome Beau\\
    \small Independent Researcher\\
    \small \href{mailto:jerome.beau@cosmochrony.org}{jerome.beau@cosmochrony.org}
    }
    \date{2026\\[2pt]\small Version 2.0}
  \maketitle

  \begin{abstract}
    Q10 proposes a conditional route from the O-series spectral universality hypothesis~[U]
    toward the isotropy identification $A_H = 2$; its coefficient and geometric steps
    require a separate audit:
    $\max_c |\sigma_c(n) - \sigma_*(n)| \leq \varepsilon(q)\,\sigma_*(n)$
    uniformly over the fitting window $n \leq n_*(q)$.
    Version~1 of this paper claimed a proof of~[U] with rate $O(q^{-1/2})$, resting on a
    Lipschitz continuity of the Weil fingerprint energy in the reduced character
    $\theta = c/q$, itself derived from an operator perturbation estimate on the
    Heisenberg\textendash Schr\"odinger generators.
    That derivation is false, and this version withdraws the claim.
    We prove instead an exact obstruction: for an odd prime $q$, the multiplication
    generators of the finite Heisenberg\textendash Schr\"odinger representations at any two
    distinct central characters are
    \emph{equidistant}, at operator distance exactly $2\cos(\pi/2q)$, independently of how
    close the characters are.
    The generator route therefore supplies no modulus of continuity in $\theta$ whatever,
    not merely a weaker one, and the equicontinuity input of the former proof is
    unavailable.
    Two further gaps are recorded and retyped rather than repaired: the fingerprint energy
    of O25 is an integer Gram\textendash Schmidt rank increment on a three-component block,
    which version~1 silently replaced by a compressed operator norm; and the
    $O(q^{-1/2})$ Gromov\textendash Hausdorff rate attributed to Q5b appears in neither
    cited theorem, both of which are rate-free.
    Hypothesis~[U], its rate, and the identification $A_H = 2$ are open on the inputs
    available here; nothing below shows~[U] to be false.
    \textbf{Interpretive outlook.}
    The obstruction is informative rather than merely negative: it locates universality, if
    it holds, in the rank increment of the fingerprint filtration rather than in any norm
    continuity of the representations, and it states the two estimates a future proof must
    supply.
  \end{abstract}

  \bigskip
  \noindent\textbf{Keywords.}
  Weil fingerprint; Heisenberg\textendash Schr\"odinger representation; central character;
  operator perturbation;
  spectral universality; equicontinuity; Gram\textendash Schmidt rank increment;
  withdrawn claim; Cosmochrony.

  \noindent\textbf{MSC 2020.}
  22E25, 22E50, 47A55, 53C23, 81R05.

  \clearpage
  \tableofcontents
  \clearpage

  \section{Introduction}

    \subsection{Context}

      Paper Q10~\cite{Beau2026q10} of the Cosmochrony spectral geometry programme proposes
      the following conditional chain, whose coefficient and geometric implications are
      not adjudicated here:
      \begin{equation}\label{eq:isotropy-chain}
        \text{[U]} \;\Longrightarrow\; A_H(q) \to 2 \;\Longrightarrow\;
        g^{\mu\nu} = \mathrm{diag}(-A_\tau, 2, 2, 2).
      \end{equation}
      The universality hypothesis~[U] (Q10 Definition~3.1) asserts the existence of a universal
      power-law profile $\sigma_*(n)$ such that
      \begin{equation}\label{eq:U}
        \max_{c \in (\mathbb{Z}/q\mathbb{Z})^\times}
        |\sigma_c(n) - \sigma_*(n)| \;\leq\; \varepsilon(q)\,\sigma_*(n),
      \end{equation}
      uniformly for $n \leq n_*(q)$, with $\varepsilon(q) \to 0$.
      Q10 Proposition~3.2 provides structural support for~[U] and numerical evidence
      (O25 inter-pair variance), but leaves the proof of~\eqref{eq:U} open.

      Version~1 of the present paper claimed to close this gap.
      It did not.
      This version states what the failed argument establishes, what it does not, and what a
      proof would still have to supply.

    \subsection{What changes in version 2.0}

      Three independent defects are recorded.

      \begin{enumerate}
        \item \textbf{The generator estimate is false.}
        Version~1 asserted
        $\|\rho_{q,c}(s) - \rho_{q,c'}(s)\|_{\mathrm{op}} \leq 2\pi|c-c'|/q$.
        The true value is computed in Lemma~\ref{lem:generator-perturb}: for the
        multiplication generator it equals $2\cos(\pi/2q)$ for \emph{every} pair of distinct
        characters. At $q = 61$ and $c' = c+1$ the asserted bound is $0.1030$ and the true
        distance is $1.9993$.
        \item \textbf{The observable was silently replaced.}
        O25~\cite{Beau2026a29} defines $\delta r_n$ as the number of shell-$n$ Weil
        fingerprint vectors linearly independent of the Gram\textendash Schmidt span of the
        earlier shells, measured on a three-component block $(c_1,c_2,c_3)$ under a sampling
        protocol. Version~1 set this integer equal to a compressed operator norm
        $\|\Pi_{S_n}\,d\rhoc\,\Pi_{S_n}\|_{\mathrm{op}}$ without any identification, and
        without defining $\Pi_{S_n}$ as a subspace of $\Cq$.
        \item \textbf{The imported rate does not exist.}
        Version~1 attributed a Gromov\textendash Hausdorff rate $O(q^{-1/2})$ to Q5b
        Theorem~2.1. That theorem is the Bass\textendash Guivarc'h ball growth
        $|B_n| \sim Cn^4$; Q5b's Carnot convergence theorem is qualitative and states no rate.
      \end{enumerate}

      Defect~1 is fatal by itself and is not a matter of constants: the obstruction below
      shows the route yields no modulus of continuity at all.

    \subsection{Main result}

      \begin{theorem}[Equidistance obstruction to the generator route]\label{thm:main}
        Let $q \geq 5$ be prime.
        For all $c \neq c'$ in $(\mathbb{Z}/q\mathbb{Z})^\times$, the Heisenberg multiplication
        generators satisfy
        \begin{equation}\label{eq:rate}
          \bigl\|\rho_{q,c}(X) - \rho_{q,c'}(X)\bigr\|_{\mathrm{op}}
          \;=\; 2\cos\!\Bigl(\frac{\pi}{2q}\Bigr) \;\geq\; 2\cos\frac{\pi}{10} > 1.90,
        \end{equation}
        a value independent of $c$ and $c'$ and tending to $2$ as $q \to \infty$.
        Consequently there is no constant $K$, independent of $q$, for which
        $\|\rho_{q,c}(X) - \rho_{q,c'}(X)\|_{\mathrm{op}} \leq K|c-c'|/q$ holds for all
        $q$ and all $c \neq c'$; and no modulus of continuity in the reduced character
        $\theta = c/q$ can be extracted from the operator distance between the generators.
        In particular the equicontinuity input used by version~1 is unavailable, and
        hypothesis~[U] is not established by that route.
      \end{theorem}

      Theorem~\ref{thm:main} is proved in Section~\ref{sec:proof} from the exact computation
      of Lemma~\ref{lem:generator-perturb}.
      It is a statement about their Heisenberg multiplication generators, not about~[U]:
      it closes one route and
      leaves the hypothesis open.

    \subsection{Structure of the argument}

      Section~\ref{sec:setup} fixes the observable and separates the quantity O25 measures
      from the quantity version~1 estimated.
      Section~\ref{sec:lipschitz} computes the exact generator distance and retypes the
      withdrawn Lipschitz lemma.
      Sections~\ref{sec:largetheta} and~\ref{sec:smalltheta} retype the two regimes of the
      former proof, each of which fails independently of the generator estimate.
      Section~\ref{sec:proof} proves Theorem~\ref{thm:main}.
      Section~\ref{sec:consequences} withdraws the downstream consequences and states the
      estimates a proof of~[U] would require.

  \section{Setup and Notation}\label{sec:setup}

    \subsection{Weil fingerprint energy}

      Fix a prime $q \geq 5$ and a character $c \in (\mathbb{Z}/q\mathbb{Z})^\times$.
      The finite Heisenberg\textendash Schr\"odinger representation of
      $G_q = \Heis_3(\mathbb{Z}/q\mathbb{Z})$ at central character
      $e^{2\pi i c/q}$ is denoted $\rhoc : G_q \to U(\Cq)$, acting on $\Cq = \mathbb{C}^q$ by
      \begin{equation}\label{eq:weil-action}
        (\rhoc(X)f)(k) = e^{2\pi i ck/q}\,f(k), \qquad
        (\rhoc(Y)f)(k) = f(k+1), \qquad k \in \mathbb{Z}/q\mathbb{Z}.
      \end{equation}
      Its associated Weil action implements symplectic automorphisms and is a distinct
      representation; the multiplication and shift operators in~\eqref{eq:weil-action}
      belong to the Heisenberg action.
      The observable of the O-series campaign is
      \begin{equation}\label{eq:sigc}
        \sigma_{c,B}(n) := \frac{\delta r_n(c,B)}{|S_n|},
      \end{equation}
      where $B=(c,c_2,c_3)$ is the sampled three-component block and, following
      O25~\cite{Beau2026a29}, $\delta r_n(c,B)$ is the \emph{number} of Weil
      fingerprint vectors in BFS shell~$S_n$ that are linearly independent of the
      Gram\textendash Schmidt span built over shells $0,\dots,n-1$, and $|S_n|$ is the $n$-th
      shell size of $\mathrm{Cay}(G_q, S_q)$.
      O25 measures this under a block-sampling protocol; it is not presented there as a
      deterministic function of a single central character.
      The historical notation
      $\sigc(n)$ in~\eqref{eq:U} suppresses the block and requires a specified aggregation
      before~[U] can be tested against the O25 measurement.
      The universal profile posited by~[U] is $\sigma_*(n) = C_*\,n^{-\dpair/2}$ with
      $\dpair \approx 7.44$.

      Two properties of~\eqref{eq:sigc} are used below and were not respected by version~1.
      First, $\delta r_n(c,B)$ is an integer in $\{0,1,\dots,|S_n|\}$, so
      $\sigma_{c,B}(n) \in [0,1]$
      takes finitely many values for each $q$ and $n$.
      Second, the BFS shells $S_n$ are subsets of the group $G_q$, of order $q^3$, and carry
      no canonical structure of subspace of the $q$-dimensional carrier $\Cq$.

    \subsection{Reduced character and parity scope}

      Define the reduced character $\theta := c/q \in (0,1)$.
      The former argument used the parity map $c\leftrightarrow q-c$ to restrict attention
      to $\theta\in(0,1/2]$.
      O22 does not prove equality of the O25 rank increments under this map, and parity
      covariance of a different response family does not supply that equality either.
      No parity reduction for the sampled-block observable is imported here.

    \subsection{What Q5b supplies}

      Version~1 asserted a Gromov\textendash Hausdorff rate
      \begin{equation}\label{eq:bfs-rate}
        d_{\mathrm{GH}}\!\bigl(\mathrm{BFS}_n(G_q),\, B_n^{CC}\bigr) \;=\; O(q^{-1/2}),
      \end{equation}
      attributing it to Q5b Theorem~2.1~\cite{Beau2026q5b}.
      Equation~\eqref{eq:bfs-rate} is \emph{not} a Q5b result and is not used below.
      Q5b Theorem~2.1 is the Bass\textendash Guivarc'h growth law $|B_n| \sim Cn^4$ for the
      discrete Heisenberg group, from which the homogeneous dimension $D_{\mathrm{hom}} = 4$ is read.
      Q5b's Carnot convergence theorem states Gromov\textendash Hausdorff convergence of the
      rescaled BFS metric space to a Carnot\textendash Carath\'eodory ball, qualitatively,
      with no rate and with no estimate transferring a metric rate to the
      Gram\textendash Schmidt independence count of~\eqref{eq:sigc}.
      Supplying such a transfer is one of the two missing estimates listed in
      Remark~\ref{rem:route}.

  \section{The Generator Route and Its Failure}\label{sec:lipschitz}

    \begin{lemma}[Heisenberg generator distance: exact value]\label{lem:generator-perturb}
      Let $q \geq 5$ be prime and $c, c' \in (\mathbb{Z}/q\mathbb{Z})^\times$ with
      $c \neq c'$. In the coordinates of~\eqref{eq:weil-action},
      \begin{equation}\label{eq:gen-perturb}
        \bigl\|\rho_{q,c}(X) - \rho_{q,c'}(X)\bigr\|_{\mathrm{op}}
        \;=\; 2\max_{0 \leq k < q}\Bigl|\sin\frac{\pi(c-c')k}{q}\Bigr|
        \;=\; 2\cos\frac{\pi}{2q},
      \end{equation}
      independently of $c$ and $c'$, while $\rho_{q,c}(Y) = \rho_{q,c'}(Y)$.
      The same value holds for $X^{-1}$ by unitarity.
    \end{lemma}

    \begin{proof}
      Write $d := c - c' \not\equiv 0$.
      By~\eqref{eq:weil-action}, $\rho_{q,c}(X) - \rho_{q,c'}(X)$ is diagonal, with $k$-th entry
      \begin{equation*}
        e^{2\pi i ck/q} - e^{2\pi i c'k/q}
        \;=\; e^{2\pi i c'k/q}\bigl(e^{2\pi i dk/q} - 1\bigr).
      \end{equation*}
      Hence that entry has modulus
      $|e^{2\pi i dk/q} - 1| = 2\,|\sin(\pi dk/q)|$, and the operator norm of a diagonal
      matrix is the largest modulus of its entries, which gives the first equality.

      Since $q$ is prime and $d \not\equiv 0$, the map $k \mapsto dk$ permutes
      $\mathbb{Z}/q\mathbb{Z}$, so
      $\max_k |\sin(\pi dk/q)| = \max_{0 \leq m < q} |\sin(\pi m/q)|$, a quantity independent
      of $d$.
      For odd $q$ the maximum is attained at $m = (q \pm 1)/2$ and equals
      $\sin\bigl(\pi(q-1)/2q\bigr) = \cos(\pi/2q)$.
      The translation $\rho_{q,c}(Y)$ does not depend on $c$, giving the last claim.
    \end{proof}

    \begin{lemma}[Lipschitz continuity of $\sigma_c(n)$: withdrawn]\label{lem:lipschitz}
      Version~1 asserted the existence of a constant $L > 0$, independent of $q$ and
      $n \leq n_*(q)$, with
      \begin{equation}\label{eq:lipschitz}
        |\sigma_c(n) - \sigma_{c'}(n)| \;\leq\; L\,\frac{|c - c'|}{q}
        \qquad \text{for all } c, c' \in (\mathbb{Z}/q\mathbb{Z})^\times,
      \end{equation}
      and hence that $\theta \mapsto \sigma_{\lfloor\theta q\rfloor}(n)$ is $L$-Lipschitz on
      $(0,1/2]$ uniformly in $q$ and $n$.
      \textbf{Statement~\eqref{eq:lipschitz} is not established}, and the argument given for
      it fails at two independent points.
    \end{lemma}

    \begin{proof}[Discussion]
      The former derivation set
      \begin{equation*}
        \delta r_n(c) \;=\; \|\Pi_{S_n}\,d\rhoc\,\Pi_{S_n}\|_{\mathrm{op}},
        \qquad d\rhoc \;=\; \sum_{s \in S_q}\bigl(\rhoc(s) - I\bigr),
      \end{equation*}
      and then combined the operator-norm stability inequality
      \begin{equation}\label{eq:op-stability}
        \bigl|\|PAP\|_{\mathrm{op}} - \|PBP\|_{\mathrm{op}}\bigr|
        \;\leq\; \|P(A - B)P\|_{\mathrm{op}} \;\leq\; \|A - B\|_{\mathrm{op}}
      \end{equation}
      with the asserted generator bound $2\pi|c-c'|/q$, reaching $L = 8\pi$.

      Inequality~\eqref{eq:op-stability} is correct for bounded operators and a projection
      $P$, and is retained.
      Both of its inputs fail.

      \emph{First}, the claimed identification of $\delta r_n(c)$ with a compressed operator norm is
      not proved anywhere, and the two quantities are of different kinds: by
      Section~\ref{sec:setup}, the actual $\delta r_n(c,B)$ is an integer count of newly independent
      fingerprint vectors on a three-component block, whereas the right-hand side is a
      continuous function of one central character.
      The projection $\Pi_{S_n}$ is moreover never defined: BFS shells are subsets of $G_q$,
      not subspaces of $\Cq$.
      Whatever~\eqref{eq:op-stability} bounds, it is not the O25 observable.

      \emph{Second}, Lemma~\ref{lem:generator-perturb} invalidates the former
      generator-by-generator triangle bound, but alone says nothing about cancellation in
      their sum: the $X+X^{-1}$ terms even coincide when $c'=-c$.
      A separate adjacent-character example does rule out the claimed modulus for this sum.
      For $c=1$, $c'=2$, and $k=(q-1)/2$, the $Y$ and identity terms cancel, while the
      difference of the $X+X^{-1}$ diagonals has absolute value
      $2\cos(\pi/q)+2\cos(2\pi/q)\to4$ as $q\to\infty$.
      Thus $\|d\rho_{q,1}-d\rho_{q,2}\|_{\mathrm{op}}$ stays bounded away from zero
      although $|1-2|/q\to0$.
      The former proof obtained its factor by writing
      $\max_{0 \leq k < q} 2\pi|c-c'|\,|k|/q \leq 2\pi\,\frac{|c-c'|}{q}\cdot\frac{q-1}{q}$;
      the left-hand side equals $2\pi|c-c'|(q-1)/q$, larger by a factor $q-1$.
    \end{proof}

    \begin{remark}[The observable is quantised]\label{rem:uniform-n}
      Version~1 discussed the conservative step $|S_n| \geq 1$ in the divisor
      of~\eqref{eq:sigc}.
      The binding difficulty lies elsewhere.
      For any fixed sampled blocks, $\delta r_n(c,B)$ is an integer and
      $\sigma_{c,B}(n)$ takes values in the finite grid
      $\{0, 1/|S_n|, \dots, 1\}$; two block measurements either agree or differ by at
      least $1/|S_n|$.
      A statement such as~\eqref{eq:lipschitz} with $|c-c'| = 1$ therefore forces
      equality of the compared block measurements as soon as $L|S_n| < q$, which is a rigidity
      claim about the fingerprint filtration, not a continuity estimate.
      Any future Lipschitz-type input must be stated for the quantised observable, or for an
      explicitly constructed continuous surrogate together with a proved comparison.
    \end{remark}

    \begin{remark}[Absolute and relative error]\label{rem:absolute-relative}
      Hypothesis~\eqref{eq:U} is a \emph{relative} bound: the error is measured against
      $\sigma_*(n) = C_*n^{-\dpair/2}$, which decreases in $n$ across the fitting window.
      Every estimate asserted in version~1 downstream of Lemma~\ref{lem:lipschitz} was
      \emph{absolute}, of the form $\leq Cq^{-1/2}$, and the two were exchanged without
      comment in the final combination.
      An absolute bound does not imply the relative one, and at the upper end $n = n_*(q)$ of
      the window the gap is the whole point of~[U].
      A proof must either establish the relative form directly or supply a separate lower
      bound on $\sigma_*(n)$ over the window.
    \end{remark}

  \section{Large-$\theta$ Regime: What the Former Argument Assumed}\label{sec:largetheta}

    \begin{proposition}[Pointwise convergence at fixed $\theta$: not established]\label{prop:pointwise}
      Version~1 asserted that for each fixed $\theta \in (0,1/2]$ and
      $c = c(q) = \lfloor\theta q\rfloor$,
      \begin{equation}\label{eq:pointwise}
        \sigma_{c(q)}(n) \;\to\; \sigma_*(n) \quad \text{as } q \to \infty,
        \qquad \text{with } |\sigma_{c(q)}(n) - \sigma_*(n)| = O(q^{-1/2})
      \end{equation}
      uniformly for $n \leq n_*(q)$.
      Neither the convergence nor the rate is established.
    \end{proposition}

    \begin{proof}[Discussion]
      The former three-step argument fails at each step.

      \textbf{Step~1} invoked an operator-norm convergence of
      $\rho_{q,\lfloor\theta q\rfloor}$ to a Schr\"odinger representation $\pi_\theta$,
      attributed to ``Q5a~T2''.
      Q5a version~3.2~\cite{Beau2026q5a} has withdrawn the earlier spatial continuum
      derivation and records a zero-form theorem together with a normalisation no-go; it does
      not supply the cited convergence.
      A current and exact source must be identified before this step can be used.

      \textbf{Step~2} transported the convergence to $\sigma_c(n)$ through the identification
      retyped in Lemma~\ref{lem:lipschitz}, and through the rate~\eqref{eq:bfs-rate} that
      Q5b does not provide.

      \textbf{Step~3} concluded that the limit is independent of $\theta$ because
      $U_\theta \pi_\theta(g) U_\theta^{-1} = \pi_1(\delta_\theta(g))$ for the Carnot
      dilation $\delta_\theta$.
      This displayed relation changes the group argument: it is an equivalence after pullback
      along $\delta_\theta$, not a unitary equivalence of two representations of the same
      fixed group, and the central character is not preserved by it.
      Neither the BFS generator at fixed scale nor the fingerprint observable was shown
      invariant under this dilation.
    \end{proof}

    \begin{proposition}[Uniform convergence on {$[\theta_1,\,1/2]$}: not established]\label{prop:uniform-large}
      Version~1 asserted, for fixed $\theta_1 \in (0,1/4)$,
      \begin{equation}\label{eq:uniform-large}
        \max_{c\,:\,c/q \in [\theta_1, 1/2]}
        |\sigma_c(n) - \sigma_*(n)| \;=\; O(q^{-1/2})
      \end{equation}
      uniformly for $n \leq n_*(q)$.
      This does not follow from the stated inputs.
    \end{proposition}

    \begin{proof}[Discussion]
      The argument applied the Arzel\`a\textendash Ascoli theorem to the family of maps
      $\theta \mapsto \sigma_{\lfloor\theta q\rfloor}(n)$, indexed by $q$ and $n$,
      using Lemma~\ref{lem:lipschitz} for equicontinuity and
      Proposition~\ref{prop:pointwise} for pointwise convergence.
      Both inputs are withdrawn above.

      Two further points survive their withdrawal and constrain any replacement.
      Equicontinuity together with pointwise convergence yields uniform convergence on a
      compact set, but yields \emph{no rate}: a rate requires pointwise constants uniform in
      $\theta$ and $n$, which version~1 did not state.
      And the step function $\theta \mapsto \sigma_{\lfloor\theta q\rfloor}(n)$ is not
      Lipschitz across grid jumps unless adjacent values coincide, so equicontinuity of the
      family is a statement about the grid values, not about the interpolant.
    \end{proof}

  \section{Small-$\theta$ Regime}\label{sec:smalltheta}

    \begin{lemma}[Uniform bound for small characters: not established]\label{lem:smalltheta}
      Version~1 asserted, for fixed $\theta_1 \in (0,1/4)$ and all $\theta \in (0,\theta_1)$,
      \begin{equation}\label{eq:smalltheta}
        |\sigma_c(n) - \sigma_*(n)| \;\leq\; C'\,q^{-1/2}
      \end{equation}
      for a constant $C'$ depending only on $L$, $\theta_1$ and the large-$\theta$ bound.
      The claim is withdrawn.
    \end{lemma}

    \begin{proof}[Discussion]
      Beyond its reliance on the withdrawn Lemma~\ref{lem:lipschitz} and
      Proposition~\ref{prop:uniform-large}, the argument contains two errors of its own.

      The triangle inequality
      \begin{equation}\label{eq:triangle}
        |\sigma_c(n) - \sigma_*(n)|
        \;\leq\; |\sigma_c(n) - \sigma_{c_0}(n)| + |\sigma_{c_0}(n) - \sigma_*(n)|,
        \qquad c_0 = \lfloor\theta_1 q\rfloor,
      \end{equation}
      gives a first term bounded by $L\theta_1$, a constant, and a second bounded by
      $Bq^{-1/2}$.
      Version~1 then stated that the constant term ``is dominated'' once
      $L\theta_1 \leq Bq^{-1/2}$, which reverses the inequality: for large $q$ the right-hand
      side tends to zero and the condition fails, rather than holding eventually.

      Version~1 next set $\theta_1 = q^{-1/2}$ to make the constant term $O(q^{-1/2})$.
      But $\theta_1$ is fixed in the hypothesis of Proposition~\ref{prop:uniform-large}, and
      its constants $B$ depend on that fixed $\theta_1$; substituting a $q$-dependent
      $\theta_1$ afterwards does not preserve them.
      A shrinking interval requires its own estimate.
    \end{proof}

    \begin{remark}[The small-$\theta$ regime is not free]\label{rem:no-BI}
      Version~1 concluded that the small-$\theta$ regime needs no Born\textendash Infeld
      input, being absorbed by $\theta_1 = q^{-1/2}$.
      With Lemma~\ref{lem:smalltheta} withdrawn, no such absorption is available, and the
      regime $\theta \to 0$ is the one where the fingerprint filtration is least controlled.
      Whether a Born\textendash Infeld input is needed there is open; the present paper
      neither requires nor excludes one.
    \end{remark}

  \section{Proof of Theorem~\ref{thm:main}}\label{sec:proof}

    \begin{proof}[Proof of Theorem~\ref{thm:main}]
      The equality in~\eqref{eq:rate} is Lemma~\ref{lem:generator-perturb}.
      The stated lower bound follows since $q \mapsto \cos(\pi/2q)$ increases in $q$, so for
      $q \geq 5$ one has $2\cos(\pi/2q) \geq 2\cos(\pi/10) = 1.9021\ldots > 1.90$, and
      $2\cos(\pi/2q) \to 2$ as $q \to \infty$.

      Suppose a constant $K$, independent of $q$, satisfied
      $\|\rho_{q,c}(X) - \rho_{q,c'}(X)\|_{\mathrm{op}} \leq K|c-c'|/q$ for all $q$ and all
      $c \neq c'$.
      Taking $c' = c+1$ gives $2\cos(\pi/2q) \leq K/q$, hence $K \geq 2q\cos(\pi/2q)$, which
      is unbounded in $q$.
      No such $K$ exists.

      For the last statement, a modulus of continuity extracted from the generator distance
      would have to make $\|\rho_{q,c}(X) - \rho_{q,c'}(X)\|_{\mathrm{op}}$ small when
      $|\theta - \theta'| = |c-c'|/q$ is small.
      By~\eqref{eq:gen-perturb} this distance does not depend on $|c - c'|$ at all and is
      bounded below by $1.90$ uniformly, so no function $\omega$ with
      $\omega(0^+) = 0$ can dominate it.
      Since the equicontinuity used in Proposition~\ref{prop:uniform-large} was derived
      solely from this distance through~\eqref{eq:op-stability}, that input is unavailable.
    \end{proof}

    \begin{remark}[Numerical illustration]\label{rem:numerics}
      For $c' = c+1$ the asserted bound $2\pi/q$ of version~1 and the true distance
      $2\cos(\pi/2q)$ diverge as $q$ grows: at $q = 61$ they are $0.1030$ and $1.9993$; at
      $q = 211$, $0.0298$ and $1.9999$; at $q = 1009$, $0.0062$ and $2.0000$.
      The ratio grows linearly in $q$, which is the factor $q-1$ lost in the former
      maximisation.
    \end{remark}

  \section{Consequences and Open Estimates}\label{sec:consequences}

    \begin{corollary}[Closure of Q10: withdrawn]\label{cor:q10}
      Version~1 asserted that hypothesis~[U] of Q10~\cite{Beau2026q10} holds within the
      Q5a\textendash Q5b framework, and drew three consequences.
      All are withdrawn:
      \begin{enumerate}
        \item[(i)] $A_H(q) \to 2$ with rate $O(q^{-1/2})$ is \emph{not} established here;
        Q10's proposed coefficient step also awaits separate audit and remains conditional
        on~[U], which remains open.
        \item[(ii)] The effective co-metric $g^{\mu\nu} = \mathrm{diag}(-A_\tau,2,2,2)$ is not
        established by this paper.
        \item[(iii)] The claim that the Bridge conjecture of Q7~\cite{Beau2026q7} is resolved
        in the isotropic case is withdrawn on two counts: it depended on~(i), and Q7
        version~2.0 records the bridge as a missing identification on a supplied carrier,
        not as a conjecture awaiting only an isotropy input.
      \end{enumerate}
      Nothing here shows~[U] to be false; the hypothesis is open, as are its consequences.
    \end{corollary}

    \begin{remark}[Rate versus empirical exponent]\label{rem:rate}
      Version~1 compared its analytical rate $O(q^{-1/2})$ with the empirical exponent
      $\approx q^{-0.44}$ of the O25 campaign and read the proximity as corroboration.
      With no analytical rate proved, the comparison has nothing to corroborate: a fitted
      exponent cannot validate a theorem, and agreement between a withdrawn rate and a fitted
      one is not evidence for either.
      The O25 numbers retain their own status as measurements of the fingerprint observable
      at $q \leq 601$.
    \end{remark}

    \begin{remark}[What is proved and what is assumed]\label{rem:epistemic}
      This version proves Theorem~\ref{thm:main} and Lemma~\ref{lem:generator-perturb}, which
      are unconditional statements about finite Heisenberg\textendash Schr\"odinger
      representations at distinct central characters, and retains
      inequality~\eqref{eq:op-stability} at its own scope.
      It proves nothing about~[U].
      The Q5a\textendash Q5b hypotheses are not used, so the paper no longer stands or falls
      with them.
      Version~1 described its conditional level as that of Q9 and Q10; that comparison is
      also withdrawn, since Q9 version~2.0 proves a conditional kinetic-limit theorem and
      explicitly does not discharge the lifting hypothesis~[H-lift]~\cite{Beau2026q9}.
    \end{remark}

    \begin{remark}[The two estimates a proof would require]\label{rem:route}
      A proof of~[U] along any route resembling the former one must supply two things that no
      current paper of the programme provides.
      \emph{An observable estimate}: a proved comparison between the O25
      Gram\textendash Schmidt rank increment of~\eqref{eq:sigc}, on its three-component
      block protocol, and whatever analytic quantity is to be estimated, including a
      definition of the shell compression on $\Cq$.
      \emph{A transfer estimate}: a quantitative passage from metric convergence of BFS balls
      to the independence count, with constants uniform in $n$ across the fitting window, and
      stated in the relative form demanded by~\eqref{eq:U} rather than an absolute one
      (Remark~\ref{rem:absolute-relative}).
      \textbf{Interpretive outlook.}
      Theorem~\ref{thm:main} suggests where not to look.
      The Heisenberg multiplication generators at distinct central characters are
      mutually as far apart as
      unitaries can almost be, uniformly in $q$; if the fingerprint energies nevertheless
      converge to a common profile, that universality is a property of the rank increment of
      the filtration, invisible to the operator distance between the representations that
      generate it.
      This reading is an interpretation of the obstruction, not a consequence of it.
    \end{remark}

  \section*{Acknowledgements}

    The author thanks the Cosmochrony programme for the O25 numerical data.
    The defects recorded in this version were identified in an independent audit of
    version~1.

    \clearpage
      \bibliography{cosmochrony-bibliography,references}

\end{document}
